\exercisesection{非典范分解的使用}

\begin{enumerate}
\def\labelenumi{\arabic{enumi})}
\tightlist
\item
  设 \(\rho: A \to B\) 是环同态，M, N 为两个左 B-模，P 为右 B-模。
\end{enumerate}

\begin{enumerate}
\def\labelenumi{\alph{enumi})}
\tightlist
\item
  设 \(L_{B}\) （相应地， \(L_{A}\) ）是 B-模（相应地，A-模）M 的一个投射分解。证明存在一个 A-线性的分解态射 \(L_{A} \rightarrow L_{B}\) ，在同伦意义下唯一；由此推出零次的分次 k-同态
\end{enumerate}

\[
\alpha : \operatorname{Tor} ^ {\mathrm{A}} (\mathrm{P}, \mathrm{M}) \rightarrow \operatorname{Tor} ^ {\mathrm{B}} (\mathrm{P}, \mathrm{M}) \quad \beta : \operatorname{Ext} _ {\mathrm{B}} (\mathrm{M}, \mathrm{N}) \rightarrow \operatorname{Ext} _ {\mathrm{A}} (\mathrm{M}, \mathrm{N}).
\]

\begin{enumerate}
\def\labelenumi{\alph{enumi})}
\setcounter{enumi}{1}
\tightlist
\item
  考虑典范同态（X，第 109 页）
\end{enumerate}

\[
\begin{array}{l l} \lambda : \operatorname{Tor} ^ {\mathsf {A}} (\mathsf {P}, \mathsf {M}) \to \operatorname{Tor} ^ {\mathsf {B}} (\mathsf {P} \otimes_ {\mathsf {A}} \mathsf {B}, \mathsf {M}) & \mu : \operatorname{Tor} ^ {\mathsf {A}} (\mathsf {P}, \mathsf {M}) \to \operatorname{Tor} ^ {\mathsf {B}} (\mathsf {P}, \mathsf {B} \otimes_ {\mathsf {A}} \mathsf {M}) \\ v: \operatorname{Ext} _ {\mathsf {B}} (\mathsf {B} \otimes_ {\mathsf {A}} \mathsf {M}, \mathsf {N}) \to \operatorname{Ext} _ {\mathsf {A}} (\mathsf {M}, \mathsf {N}) & \rho : \operatorname{Ext} _ {\mathsf {B}} (\mathsf {M}, \operatorname{Hom} _ {\mathsf {A}} (\mathsf {B}, \mathsf {N})) \to \operatorname{Ext} _ {\mathsf {A}} (\mathsf {M}, \mathsf {N}). \end{array}
\]

设 \(m: \mathbf{B} \otimes_{\mathbf{A}} \mathbf{M} \to \mathbf{M}\) , \(p: \mathbf{P} \otimes_{\mathbf{A}} \mathbf{B} \to \mathbf{B}\) 和 \(n: \mathbf{N} \to \operatorname{Hom}_{\mathbf{A}}(\mathbf{B}, \mathbf{N})\) 为由 \(m(b \otimes x) = bx\) , \(p(y \otimes b) = yb\) 和 \(n(z)(b) = bz\) （对 \(b \in \mathbf{B}\) , \(x \in \mathbf{M}\) , \(y \in \mathbf{P}\) , \(z \in \mathbf{N}\) ）定义的 B-同态。证明下列等式

\[
\begin{array}{l} \alpha = \operatorname{Tor} (1 _ {\mathrm{P}}, m) \circ \mu = \operatorname{Tor} (p, 1 _ {\mathrm{M}}) \circ \lambda \\ \beta = v \circ \operatorname{Ext} (m, 1 _ {\mathrm{N}}) = \rho \circ \operatorname{Ext} (1 _ {\mathrm{M}}, n). \end{array}
\]

\begin{enumerate}
\def\labelenumi{\arabic{enumi})}
\setcounter{enumi}{1}
\tightlist
\item
  设 A, B 为两个 k-代数（结合的且含单位元）；设 M 为左 A-模，P 为右 B-模，N 为 (A, B)-双模。考虑同态
\end{enumerate}

\[
\sigma : \operatorname{Hom} _ {\mathrm{B}} (\mathrm{N}, \mathrm{P}) \otimes_ {\mathrm{A}} \mathrm{M} \rightarrow \operatorname{Hom} _ {\mathrm{B}} (\operatorname{Hom} _ {\mathrm{A}} (\mathrm{M}, \mathrm{N}), \mathrm{P})
\]

使得 \(\sigma(u \otimes m)(v) = u \circ v(m)\) （对 \(u \in \mathrm{Hom}_{\mathbf{B}}(\mathbf{N}, \mathbf{P})\) , \(v \in \mathrm{Hom}_{\mathbf{A}}(\mathbf{M}, \mathbf{N})\) , \(m \in \mathbf{M}\) ）（II，第 190 页，习题 6）。a) 用投射分解代替 M，借助 \(\sigma\) 定义一个 \(k\) -同态

\[
\tau : \operatorname{Tor} ^ {\mathbf {A}} \left(\operatorname{Hom} _ {\mathbf {B}} (\mathrm{N}, \mathrm{P}), \mathrm{M}\right)\rightarrow \operatorname{Homgr} _ {\mathbf {B}} \left(\operatorname{Ext} _ {\mathbf {A}} (\mathrm{M}, \mathrm{N}), \mathrm{P}\right).
\]

\begin{enumerate}
\def\labelenumi{\alph{enumi})}
\setcounter{enumi}{1}
\item
  此后假设 A 是左诺特的，且 A-模 M 是有限生成的。证明若 B-模 P 是内射的， \(\tau\) 是一个同构。
\item
  证明存在两个谱序列 E 和 'E（X，第 175–177 页，习题 14 至 17），收敛到同一个分次 k-模，使得
\end{enumerate}

\[
E _ {2} ^ {p q} = \operatorname{Tor} _ {- p} ^ {A} \left(\operatorname{Ext} _ {B} ^ {q} (N, P), M\right) \quad^ {\prime} E _ {2} ^ {p q} = \operatorname{Ext} _ {B} ^ {p} \left(\operatorname{Ext} _ {A} ^ {- q} (M, N), P\right).
\]

\begin{enumerate}
\def\labelenumi{\arabic{enumi})}
\setcounter{enumi}{2}
\item
  设 I 为内射 \(k\) -模；对每个 A-模 M，记 D(M) 为右 A-模 \(\mathrm{Hom}_{k}(M, I)\) 。设 M 和 N 为两个左 A-模，P 为右 A-模。证明对每个 \(p \geqslant 0\) ， \(k\) -模 \(\mathrm{Ext}_{A}^{p}(P, D(M))\) 与 D( \(\mathrm{Tor}_{p}^{A}(P, M)\) ) 同构；若 A 是左诺特的且 M 是有限生成的 A-模，则 \(k\) -模 \(\mathrm{Tor}_{p}^{A}(D(N), M)\) 与 D( \(\mathrm{Ext}_{A}^{p}(M, N)\) ) 同构（利用第 188 页习题 2, b) 和习题 6, b)）。
\item
  对每个群 G，记 \(\varepsilon_{\mathbf{G}}: \mathbf{Z}^{(\mathbf{G})} \to \mathbf{Z}\) 为满足 \(\varepsilon_{\mathbf{G}}(e_g) = 1\) （对所有 \(g \in \mathbf{G}\) ）的环同态；令 \(\mathrm{I}_{\mathbf{G}} = \operatorname{Ker}(\varepsilon_{\mathbf{G}})\) 。
\end{enumerate}

\begin{enumerate}
\def\labelenumi{\alph{enumi})}
\item
  证明群 \(\mathbf{H}_{1}(\mathbf{G},\mathbf{Z})\) 同构于 \(\mathrm{I}_{\mathrm{G}} / \mathrm{I}_{\mathrm{G}}^{2}\) 。
\item
  设 \((g_{\iota})_{\iota \in \mathrm{I}}\) 是 G 的一个生成族；证明元素 \((e_{g_{\iota}} - 1)\) ，对于 \(\iota \in \mathrm{I}\) ，生成理想 \(\mathbf{I}_{\mathbf{G}}\) .
\item
  假设族 \((g_{t})_{t\in I}\) 是基本的（I，第 86 页）。证明 \((e_{g_t} - 1)_{t\in I}\) 是左 \(\mathbf{Z}^{(G)}\) -模 \(I_G\) 的一个基。
\end{enumerate}

\begin{enumerate}
\def\labelenumi{\arabic{enumi})}
\setcounter{enumi}{4}
\item
  设 X 是一个集合， \(G = F(X)\) 为构造在 X 上的自由群，M 是一个 \(Z^{(G)}\) -模。证明群 \(H^q(G, M)\) 和 \(H_q(G, M)\) 对 \(q \geqslant 2\) 为零（使用练习 4）。
\item
  设 G 是一个阶为 n 的有限群，M 是一个 \(\mathbf{Z}^{(G)}\) -模。证明对每个 \(q \geqslant 1\) ，Z-模 \(\mathrm{H}^{q}(\mathrm{G}, \mathrm{M})\) 和 \(\mathrm{H}_{q}(\mathrm{G}, \mathrm{M})\) 都被 n 零化（考虑 C(G, M) 到自身的同伦 h，其定义为
\end{enumerate}

\[
h f (g _ {1}, \dots , g _ {q}) = (- 1) ^ {q} \sum_ {g \in G} f (g _ {1}, \dots , g _ {q}, g)
\]

当 \(f \in \mathbf{C}^{q+1}(\mathbf{G}, \mathbf{M})\) 和 \(g, g_1, \ldots, g_q \in \mathbf{G}\) ）。若 \(\mathbf{M}\) 是一个平坦的有限生成 \(\mathbf{Z}\) -模，则群 \(\mathrm{H}^q(\mathbf{G}, \mathbf{M})\) 和 \(\mathrm{H}_q(\mathbf{G}, \mathbf{M})\) 对 \(q \geqslant 1\) 是有限的。

\begin{enumerate}
\def\labelenumi{\arabic{enumi})}
\setcounter{enumi}{6}
\tightlist
\item
  设 G 是 n 阶循环群，M 是一个 \(\mathbf{Z}^{(G)}\) -模， \(\sigma\) 是 G 的一个生成元。记 N 为 M 中由元素 \(\sum_{g\in G}e_{g}\) 的 \(\mathbf{Z}^{(G)}\) 给出的乘法。证明 Z-模 \(\mathrm{H}^{0}(\mathrm{G},\mathrm{M})\) （相应地， \(\mathrm{H}_{0}(\mathrm{G},\mathrm{M})\) ）同构于 Ker
\end{enumerate}

于 \((1 - \sigma_{\mathbf{M}})\) （相应地，Coker \((1 - \sigma_{\mathbf{M}})\) ）；对奇数的 \(q\) ，群 \(\mathrm{H}^q (\mathrm{G},\mathrm{M})\) 和 \(\mathrm{H}_{q + 1}(\mathrm{G},\mathrm{M})\) 同构于 Ker N/Im \((1 - \sigma_{\mathbf{M}})\) ；对偶数的 \(q\) \(\geqslant 2\) ， \(\mathrm{H}^q (\mathrm{G},\mathrm{M})\) 和 \(\mathrm{H}_{q - 1}(\mathrm{G},\mathrm{M})\) 同构于 Ker \((1 - \sigma_{\mathbf{M}}) / \mathrm{Im}\) N。（利用第 178 页练习 1。）

\begin{enumerate}
\def\labelenumi{\arabic{enumi})}
\setcounter{enumi}{7}
\item
  设F为自由群，R为F的正规子群，且G = F/R。证明群H₂(G, Z)同构于((F, F) ∩ R)/(F, R)（利用第179页练习3）。
\item
  \begin{enumerate}
  \def\labelenumii{\alph{enumii})}
  \tightlist
  \item
    设 G 是一个群，M 是一个 G-模，且 g 是 G 的中心中的一个元素，使得 M 的自同态 \(m \mapsto gm - m\) 是双射的。证明群 \(\mathrm{H}^{p}(\mathrm{G}, \mathrm{M})\) 和 \(\mathrm{H}_{p}(\mathrm{G}, \mathrm{M})\) 对 \(p \geqslant 1\) 为零。 b) 设 V 是域 k 上的向量空间，且 Card(k) \(\geqslant 3\) 。证明群 \(\mathrm{H}^{p}(\mathrm{GL}(\mathrm{V}), \mathrm{V})\) 和 \(\mathrm{H}_{p}(\mathrm{GL}(\mathrm{V}), \mathrm{V})\) 对 \(p \geqslant 1\) 为零。
  \end{enumerate}
\item
  设 G 和 H 是两个群， \(u: H \to G\) 一个同态，M 一个 \(\mathbf{Z}^{(G)}\) -模。对于每一个 \(q \geqslant 0\) ，我们从同态 \(\alpha\) 和 \(\beta\) （练习1中的）推导出同态
\end{enumerate}

\[
u _ {q} ^ {\mathrm{M}}: \mathrm{H} _ {q} (\mathrm{H}, \mathrm{M}) \rightarrow \mathrm{H} _ {q} (\mathrm{G}, \mathrm{M}) \quad u _ {\mathrm{M}} ^ {q}: \mathrm{H} ^ {q} (\mathrm{G}, \mathrm{M}) \rightarrow \mathrm{H} ^ {q} (\mathrm{H}, \mathrm{M})
\]

当关于 M 没有歧义时，我们将其简记为 \(u_q\) 和 \(u^q\) 。

\begin{enumerate}
\def\labelenumi{\alph{enumi})}
\tightlist
\item
  对于 \(f \in \mathbf{C}^q(\mathbf{G}, \mathbf{M})\) ，我们用 \(f_u\) 表示 \(\mathbf{C}^q(\mathbf{H}, \mathbf{M})\) 中满足 \(f_u(h_1, \ldots, h_q) = f(u(h_1), \ldots, u(h_q))\) （对 \(h_1, \ldots, h_q \in \mathbf{H}\) ）的元素。证明映射 \(f \mapsto f_u\) 是从 \(\mathbf{C}'(\mathbf{G}, \mathbf{M})\) 到 \(\mathbf{C}'(\mathbf{H}, \mathbf{M})\) 的一个复形态射，它在同调上诱导出同态 \(\bigoplus_{q} u^q\) 。类似地，定义一个从 \(\mathbf{C}'(\mathbf{H}, \mathbf{M})\)
\end{enumerate}

到 C'(G, M) 中的态射，它诱导出（ \(\oplus u_{q}\) ）。

\begin{enumerate}
\def\labelenumi{\alph{enumi})}
\setcounter{enumi}{1}
\tightlist
\item
  假设H = G，并取u为G的一个内自同构。证明同态 \(u_{q}\) 和 \(u^{q}\) 等于恒等映射（通过对q进行归纳论证）。
\end{enumerate}

\begin{enumerate}
\def\labelenumi{\arabic{enumi})}
\setcounter{enumi}{10}
\tightlist
\item
  设G是一个群，H是G的一个正规子群，M是一个G-模。由习题10，群G通过共轭作用在Z-模H \(^{q}\) (H, M)上；由于子群H的作用是平凡的，因此这些模被赋予了一个Z \(^{(G/H)}\) -模的结构。
\end{enumerate}

\begin{enumerate}
\def\labelenumi{\alph{enumi})}
\tightlist
\item
  证明存在一个谱序列E（X，第175页，习题14），收敛到 \(\bigoplus_{n} H^{n}(G, M)\) ,
\end{enumerate}

使得 \(E_{2}^{pq} = H^{p}(G/H, H^{q}(H, M))\) （考虑 X 第 188 页习题 6, d) 的谱序列，并顾及第 109 页的典范同构 (12)）。

\begin{enumerate}
\def\labelenumi{\alph{enumi})}
\setcounter{enumi}{1}
\tightlist
\item
  证明存在一个正合序列
\end{enumerate}

\[
0 \to \mathrm{H} ^ {1} (\mathrm{G} / \mathrm{H}, \mathrm{H} ^ {0} (\mathrm{H}, \mathrm{M})) \to \mathrm{H} ^ {1} (\mathrm{G}, \mathrm{M}) \to \mathrm{H} ^ {0} (\mathrm{G} / \mathrm{H}, \mathrm{H} ^ {1} (\mathrm{H}, \mathrm{M})) \to \mathrm{H} ^ {2} (\mathrm{G} / \mathrm{H}, \mathrm{H} ^ {0} (\mathrm{H}, \mathrm{M})) \to \mathrm{H} ^ {2} (\mathrm{G}, \mathrm{M})
\]

利用H'(G, M)在C'(G, M)中的描述来描述该序列的同态。c) 定义一个与a)中类似的用于同调的谱序列。

\begin{enumerate}
\def\labelenumi{\arabic{enumi})}
\setcounter{enumi}{11}
\tightlist
\item
  设G是一个群。称一个 \(\mathbf{Z}^{(G)}\) -模M是余诱导的（相应地，诱导的），如果存在一个Z-模A使得M同构于 \(\operatorname{Hom}_{\mathbf{Z}}(\mathbf{Z}^{(G)}, \mathbf{A})\) （相应地， \(\mathbf{Z}^{(G)} \otimes_{\mathbf{Z}} \mathbf{A}\) ）。
\end{enumerate}

\begin{enumerate}
\def\labelenumi{\alph{enumi})}
\item
  每一个 \(\mathbf{Z}^{(G)}\) -模都同构于一个余诱导模的子模，并且同构于一个诱导模的商模。
\item
  设 M 为一个 \(\mathbf{Z}^{(G)}\) -模。证明以下条件是等价的：
\end{enumerate}

α) M 是某个余诱导模的直和项。

β) M 关于 Z 是内射的（第170页习题19）。

γ) 存在 M 上的 G-平均（第I卷，第136页，习题8），即一个 \(\mathbf{Z}^{(G)}\) -同态

\(m: \mathbf{C}^{1}(\mathbf{G}, \mathbf{M}) \to \mathbf{M}\) （\(\mathbf{C}^{1}(\mathbf{G}, \mathbf{M})\) 上的 G-模结构由 \((gf)(x) = gf(g^{-1}x)\) （对 \(g \in \mathbf{G}, x \in \mathbf{M}\) ）定义）使得如果 \(f_{x}\) 是取值为 \(x \in \mathbf{M}\) 的常值函数，则有 \(m(f_{x}) = x\) 。

\begin{enumerate}
\def\labelenumi{\alph{enumi})}
\setcounter{enumi}{2}
\tightlist
\item
  证明以下条件是等价的：
\end{enumerate}

α) M 是某个诱导模的直和项。

β) M 相对于 Z 是投射的（习题19，第170页）。

\(\gamma)\) 存在一个 \(\mathbf{Z}\) -自同态 \(\rho\) （定义在 \(\mathbf{M}\) 上），使得对每一个 \(x \in \mathbf{M}\) ，族 \((\rho(g^{-1} x))_{g \in G}\) 具有有限支撑，并且有 \(x = \sum g\rho(g^{-1} x)\) 。

\begin{enumerate}
\def\labelenumi{\alph{enumi})}
\setcounter{enumi}{3}
\item
  设 M 是一个内射（相应地，投射） \(\mathbf{Z}^{(G)}\) -模，相对于 \(\mathbf{Z}\) 。证明 \(\mathrm{H}^q (\mathbf{G},\mathbf{M}) = 0\) （相应地， \(\mathrm{H}_q(\mathrm{G},\mathrm{M}) = 0\) ）对每个 \(q\geqslant 1\) 成立。
\item
  若 G 是有限的，证明诱导模与余诱导模（分别相对于 Z 的投射模与内射模）的概念是一致的。
\end{enumerate}

\begin{enumerate}
\def\labelenumi{\arabic{enumi})}
\setcounter{enumi}{12}
\tightlist
\item
  设 G 是一个群，H 是 G 的一个子群，M 是一个左 \(\mathbf{Z}^{(H)}\) -模，且 N 是一个右 \(\mathbf{Z}^{(H)}\) -模。我们赋予群 \(\hat{\mathbf{M}} = \operatorname{Hom}_{\mathbf{Z}^{(H)}} (\mathbf{Z}^{(G)}, \mathbf{M})\) （相应地， \(\check{\mathbf{N}} = \mathbf{N} \otimes_{\mathbf{Z}^{(H)}} \mathbf{Z}^{(G)})\) 一个左 \(\mathbf{Z}^{(G)}\) -模（相应地，右 \(\mathbf{Z}^{(G)}\) -模）结构，它是由 \(\mathbf{Z}^{(G)}\) 的右模结构导出的。
\end{enumerate}

设 \(m:\hat{\mathbf{M}}\to \mathbf{M}\) 和 \(n:\mathbf{N}\rightarrow \check{\mathbf{N}}\) 为 \(\mathbf{Z}^{(\mathrm{H})}\) -同态，使得 \(m(u) = u(1)\) （对 \(u\in \hat{\mathbf{M}}\) ）以及

\[
n (x) = x \otimes 1 \quad \text { for } \quad x \in \mathbf {N}  ,
\]

并设 \(i: H \to G\) 是典范嵌入。证明复合同态

\[
\hat {i} ^ {q}: \mathrm{H} ^ {q} (\mathrm{G}, \hat {\mathrm{M}}) \xrightarrow {i ^ {q}} \mathrm{H} ^ {q} (\mathrm{H}, \hat {\mathrm{M}}) \xrightarrow {\mathrm{H} ^ {q} (\mathrm{H} , m)} \mathrm{H} ^ {q} (\mathrm{H}, \mathrm{M})
\]

和

\[
\check {i} _ {q}: \mathrm{H} _ {q} (\mathrm{H}, \mathrm{N}) \xrightarrow {\mathrm{H} ^ {q} (\mathrm{H} , n)} \mathrm{H} _ {q} (\mathrm{H}, \check {\mathrm{N}}) \xrightarrow {i _ {q}} \mathrm{H} _ {q} (\mathrm{G}, \check {\mathrm{N}})
\]

对所有 \(q \geqslant 0\) 都是同构。

\begin{enumerate}
\def\labelenumi{\arabic{enumi})}
\setcounter{enumi}{13}
\tightlist
\item
  都是同构。设 G 是一个群，且 H 是 G 的一个有限指数子群。我们令 \(n = (G:H)\) .
\end{enumerate}

\begin{enumerate}
\def\labelenumi{\alph{enumi})}
\tightlist
\item
  设 Q, M 为两个 \(\mathbf{Z}^{(G)}\) -模。证明存在一个 \(\mathbf{Z}\) -同态
\end{enumerate}

\[
t: \operatorname{Hom} _ {\mathbf {Z} ^ {(H)}} (Q, M) \rightarrow \operatorname{Hom} _ {\mathbf {Z} ^ {(G)}} (Q, M)
\]

使得，对于 G 模 H 的类的每一组代表元系 \(g_{1}, \ldots, g_{n}\) ，

\[
(t u) (q) = \sum_ {i = 1} ^ {n} g _ {i} u \left(g _ {i} ^ {- 1} q\right) \quad \text { for } \quad q \in \mathrm{Q}, \quad u \in \operatorname{Hom} _ {\mathbf {Z} (\mathrm{H})} (\mathrm{Q}, \mathrm{M}).
\]

若 \(j: \operatorname{Hom}_{\mathbf{Z}(\mathbf{G})}(\mathbf{Q}, \mathbf{M}) \to \operatorname{Hom}_{\mathbf{Z}(\mathbf{H})}(\mathbf{Q}, \mathbf{M})\) 是典范嵌入，我们有 \(t \circ j = n\) 倍恒等映射。

\begin{enumerate}
\def\labelenumi{\alph{enumi})}
\setcounter{enumi}{1}
\tightlist
\item
  设 N 为一个右 \(\mathbf{Z}^{(G)}\) -模。使用以下记号 \(a\) )，证明存在一个 \(\mathbf{Z}\) -同态定义 \(\theta: \mathrm{N} \otimes_{\mathbf{Z}^{(G)}} \mathrm{Q} \to \mathrm{N} \otimes_{\mathbf{Z}^{(H)}} \mathrm{Q}\) 使得
\end{enumerate}

\[
\theta (n \otimes q) = \sum_ {i = 1} ^ {n} n g _ {i} ^ {- 1} \otimes g _ {i} q \quad \text { for } \quad n \in \mathbf {N}, \quad q \in \mathbf {Q}.
\]

若 \(p: \mathbf{N} \otimes_{\mathbf{Z}(\mathbf{H})} \mathbf{Q} \to \mathbf{N} \otimes_{\mathbf{Z}(\mathbf{G})} \mathbf{Q}\) 是典范映射，我们有 \(p \circ \theta = n\) 倍恒等映射。
\begin{enumerate}
\def\labelenumi{\alph{enumi})}
\setcounter{enumi}{2}
\item
  设 P 是 \(\mathbf{Z}^{(G)}\) -模 \(\mathbf{Z}\) 的一个投射分解， \(q\) 为一个整数。通过应用 \(a\) 和 \(b)\) 且取 \(\mathbf{Q} = \mathbf{P}_q\) ，定义同态 \(t^q: \mathrm{H}^q(\mathrm{H}, \mathrm{M}) \to \mathrm{H}^q(\mathrm{G}, \mathrm{M})\) 和 \(\theta_q: \mathrm{H}_q(\mathrm{G}, \mathrm{N}) \to \mathrm{H}_q(\mathrm{H}, \mathrm{N})\) 。如果 \(i\) 表示 H 到 G 的典范嵌入，我们有 \(t^q \circ i^q = n\) 倍恒等映射，以及 \(i_q \circ \theta_q = n\) 倍恒等映射。
\item
  我们考虑 \(\mathbf{Z}^{(G)}\) -模 \(\hat{\mathbf{M}} = \operatorname{Hom}_{\mathbf{Z}^{(H)}} (\mathbf{Z}^{(G)}, \mathbf{M})\) 和 \(\check{\mathbf{N}} = \mathbf{N} \otimes_{\mathbf{Z}^{(H)}} \mathbf{Z}^{(G)}\) （练习13）。取 \(\mathbf{Q} = \mathbf{Z}^{(G)}\) ，我们从同态 \(t\) 和 \(\theta\) （\(a)\) 和 \(b)\) 中的）得到同态 \(t_\mathbf{M}: \hat{\mathbf{M}} \to \mathbf{M}\) 和 \(\theta_\mathbf{N}: \mathbf{N} \to \check{\mathbf{N}}\) 。证明这些同态是 \(\mathbf{Z}^{(G)}\) -线性的。证明复合同态
\end{enumerate}

\[
\begin{array}{c} \mathrm{H} ^ {q} (\mathrm{H}, \mathrm{M}) \xrightarrow {\stackrel {{(\vec {l} _ {q}) ^ {- 1}}} {{\longrightarrow}}} \mathrm{H} ^ {q} (\mathrm{G}, \stackrel {{\hat {\mathrm{M}}}} {{\longrightarrow}}) \xrightarrow {\mathrm{H} ^ {q} (\mathrm{G} , t _ {\mathrm{M}})} \mathrm{H} ^ {q} (\mathrm{G}, \mathrm{M}) \\ (\text {resp.} \mathrm{H} _ {q} (\mathrm{G}, \mathrm{N}) \xrightarrow {\mathrm{H} _ {q} (\mathrm{G} , \theta_ {\mathrm{N}})} \mathrm{H} _ {q} (\mathrm{G}, \check {\mathrm{N}}) \xrightarrow {\stackrel {{(\vec {l} _ {q}) ^ {- 1}}} {{\longrightarrow}}} \mathrm{H} _ {q} (\mathrm{H}, \mathrm{N})) \text {is equal to } t ^ {q} (\text {resp.} \theta_ {q}). \end{array}
\]

\begin{enumerate}
\def\labelenumi{\alph{enumi})}
\setcounter{enumi}{4}
\item
  设 \((g_{1},\ldots ,g_{n})\) 是 G 关于 H 的商类的代表元系，并令 \(x\in \mathrm{H}^0 (\mathrm{H},\mathrm{M})\) 。证明我们有 \(t^0 (x) = \sum_{i = 1}^{n}g_{i}x\) 。描述 \(t^1\)
\item
  如果将群 \(\mathrm{H}_{1}(\mathrm{G},\mathbf{Z})\) 和 \(\mathrm{H}_{1}(\mathrm{H},\mathbf{Z})\) 分别与 G/(G, G) 和 H/(H, H) 等同起来，则同态 \(t_{1}\) 被等同于一个同态（“转移”）V : G/(G, G) → H/(H, H)。若 s : \(H \backslash G\) → G 是典范映射的一个截面，证明 V 是由映射 \(V_{s}\) ：G → H 通过取商得到的，其中
\end{enumerate}

\[
\mathrm{V} _ {s} (g) = \prod_ {x \in \mathrm{H} \backslash \mathrm{G}} s (x) g (s (x g)) ^ {- 1} \quad \text { for } \quad g \in \mathrm{G}.
\]

\begin{enumerate}
\def\labelenumi{\arabic{enumi})}
\setcounter{enumi}{14}
\tightlist
\item
  设G是一个群，k是一个交换域；我们赋予 k 由 G 的平凡作用导出的 \(\mathbf{Z}^{(\mathbf{G})}\) -模结构以及 \(k^{(\mathbf{G})}\) -模结构。
\end{enumerate}

\begin{enumerate}
\def\labelenumi{\alph{enumi})}
\item
  证明群 \(\mathrm{H}^q (\mathbf{G},k)\) 同构于 \(\operatorname {Ext}_{k(\mathbf{G})}^{q}(k,k)\) （对所有 \(q\geqslant 0\) ），因此具有 \(k\) -向量空间的结构；我们用 \(h^q (\mathbf{G},k)\) 表示它在 \(k\) 上的维数
\item
  假设群 G 具有由 d 个生成元和 r 个关系定义的一个表示（1，第 86 页）。证明以下不等式
\end{enumerate}

\[
d \geqslant h ^ {1} (\mathbf {G}, k) \quad \text { and } \quad r \geqslant h ^ {2} (\mathbf {G}, k)
\]

（使用习题3，d），第179页）。

¶ 16) 设 G 是一个 p-群（1，p.～72，定义 9）。我们令 \(d(\mathbf{G}) = \dim_{\mathbf{F}_{p}} \mathrm{H}^{1}(\mathbf{G}, \mathbf{F}_{p})\) 和 \(r(\mathbf{G}) = \dim_{\mathbf{F}_{p}} \mathrm{H}^{2}(\mathbf{G}, \mathbf{F}_{p})\) （参见练习 15）。

\begin{enumerate}
\def\labelenumi{\alph{enumi})}
\item
  证明 \(d(\mathbf{G})\) 是 \(\mathbf{G}\) 的生成元的最小数目（证明 \(\mathrm{H}^1 (\mathbf{G},\mathbf{F}_p)\) 同构于 \(\mathbf{G} / \mathbf{G}^p\mathbf{D}(\mathbf{G})\) 的对偶，并利用 I，第 140 页，练习 32）。
\item
  如果 Card (G) = \(p^n\) ，证明有 \(d(\mathbf{G}) \leqslant n\) 和 \(r(\mathbf{G}) \leqslant \frac{1}{2} n(n + 1)\) （通过对 \(n\) 进行归纳证明，存在 G 的一个表示，具有 \(n\) 个生成元和 \(\frac{1}{2} n(n + 1)\) 个关系）。如果 \(\mathbf{G} = (\mathbf{Z}/p\mathbf{Z})^n\) ，证明有 \(d(\mathbf{G}) = n\) 和 \(r(\mathbf{G}) = \frac{1}{2} n(n + 1)\) 。
\item
  如果 G 是非平凡的，证明不等式 \(r(\mathbf{G}) > \frac{1}{4}(d(\mathbf{G}))^{2}\) (“Golod-Shafarevich定理”：使用第182页练习15，以及VIII章第8节练习22)。
\end{enumerate}

\(d)\) 证明以下命题成立： \(r(\mathbf{G}) - d(\mathbf{G}) = \dim_{\mathbf{F}_p}(\mathbf{H}^3 (\mathbf{G},\mathbf{Z})\otimes_{\mathbf{Z}}\mathbf{F}_p).\)

\begin{enumerate}
\def\labelenumi{\arabic{enumi})}
\setcounter{enumi}{16}
\tightlist
\item
  设 G 是一个有限群。我们用 N 表示元素 \(\sum_{g\in G}e_g\) （属于 \(\mathbf{Z}^{(G)}\) ）。对每个 \(\mathbf{Z}^{(G)}\) -模 M，乘法
\end{enumerate}

由 N 通过取商定义出一个同态 \(\mathrm{N}_{0}:\mathrm{H}_{0}(\mathrm{G},\mathrm{M})\to\mathrm{H}^{0}(\mathrm{G},\mathrm{M})\) 。我们设 \(\hat{\mathrm{H}}^{0}(\mathrm{G},\mathrm{M})=\mathrm{Coker}(\mathrm{N}_{0}),\hat{\mathrm{H}}^{-1}(\mathrm{G},\mathrm{M})=\mathrm{Ker}(\mathrm{N}_{0}),\hat{\mathrm{H}}^{q}(\mathrm{G},\mathrm{M})=\mathrm{H}^{q}(\mathrm{G},\mathrm{M})\text{ 当 } q\geqslant1,\text{ 且}\)

\[
\hat {\mathrm{H}} ^ {- q} (\mathrm{G}, \mathrm{M}) = \mathrm{H} _ {q - 1} (\mathrm{G}, \mathrm{M}) \quad \text { for } \quad q \geqslant 2.
\]

\begin{enumerate}
\def\labelenumi{\alph{enumi})}
\tightlist
\item
  设 \(\mathcal{E}:0\to\mathbf{M}^{\prime}\to\mathbf{M}\to\mathbf{M}^{\prime\prime}\to0\) 为一个 \(\mathbf{Z}^{(G)}\) -模的正合序列；定义一个正合序列
\end{enumerate}

\[
\dots \rightarrow \hat {\mathrm{H}} ^ {q - 1} (\mathrm{G}, \mathrm{M} ^ {\prime \prime}) \xrightarrow {\partial^ {q - 1} (\mathrm{G} , \mathcal {E})} \hat {\mathrm{H}} ^ {q} (\mathrm{G}, \mathrm{M} ^ {\prime}) \rightarrow \hat {\mathrm{H}} ^ {q} (\mathrm{G}, \mathrm{M}) \rightarrow \hat {\mathrm{H}} ^ {q} (\mathrm{G}, \mathrm{M} ^ {\prime \prime}) \xrightarrow {\partial^ {q} (\mathrm{G} , \mathcal {E})} \hat {\mathrm{H}} ^ {q + 1} (\mathrm{G}, \mathrm{M} ^ {\prime}) \rightarrow \dots
\]

\begin{enumerate}
\def\labelenumi{\alph{enumi})}
\setcounter{enumi}{1}
\tightlist
\item
  若M相对于Z是投射的，证明有 \(\hat{\mathrm{H}}^{q}(\mathrm{G},\mathrm{M}) = 0\) 对所有 \(q\in \mathbf{Z}\) 成立（参见练习12）。c）设H是G的一个子群， \(i:\mathrm{H}\to \mathrm{G}\) 为典范嵌入。对于每一个 \(\mathbf{Z}^{(\mathrm{G})}\) -模 M 以及每个有理整数 \(q\) ，定义同态 \(\hat{i}_{\mathrm{M}}^{\mathrm{q}}:\hat{\mathrm{H}}^{\mathrm{q}}(\mathrm{G},\mathrm{M})\rightarrow \hat{\mathrm{H}}^{\mathrm{q}}(\mathrm{H},\mathrm{M})\) 和 \(\hat{t}_{\mathrm{M}}^{\mathrm{q}}:\hat{\mathrm{H}}^{\mathrm{q}}(\mathrm{H},\mathrm{M})\rightarrow \hat{\mathrm{H}}^{\mathrm{q}}(\mathrm{G},\mathrm{M})\) （若关于 M 无歧义，就简记为 \(\hat{i}^q\) 和 \(\hat{t}^q\) ）如下：对于 \(q\geqslant 1\) ，我们设 \(\hat{i}^q = i^q\) 和 \(\hat{t}^q = t^q\) （练习10和14）；对于 \(q\leqslant -2\) ，我们设 \(\hat{i}^q = \theta_{-q - 1}\) 和 \(\hat{t}^q = i_{-q - 1}\) ；最后，定义 \(\hat{i}^0\) 和 \(\hat{t}^0\) （相应地， \(\hat{i}^{-1}\) 和 \(\hat{t}^{-1}\) ）为通过取商（相应地，取子集）从 \(i^0\) 和 \(t^0\) （相应地， \(t_0\) 和 \(i_0\) ）得到的同态。证明有 \(\hat{i}^q\circ \hat{t}^q = n\) 倍恒等映射（对所有 \(q\in \mathbf{Z}\) ），其中 \(n\) 是H在G中的指数。由此推出，群 \(\hat{\mathbf{H}}^{q}(\mathbf{G},\mathbf{M})\) 被 G 的阶零化（参见第189页练习6）。
\end{enumerate}

\(d)\) 使用 \(b\) 中的记号，证明有 \(\hat{i}_{M'}^{q+1} \circ \partial^q(G, \mathcal{E}) = \partial^q(H, \mathcal{E}) \circ \hat{i}_{M''}^q\) 和

\[
\hat {t} _ {\mathbf {M} ^ {\prime}} ^ {q + 1} \circ \partial^ {q} (\mathrm{H}, \mathcal {E}) = \partial^ {q} (\mathrm{G}, \mathcal {E}) \circ \hat {t} _ {\mathbf {M} ^ {\prime \prime}} ^ {q}
\]

对所有 \(q \in \mathbf{Z}\) .

\begin{enumerate}
\def\labelenumi{\alph{enumi})}
\setcounter{enumi}{4}
\tightlist
\item
  设 \(p\) 是一个素数，H 为 G 的一个 \(p\) -Sylow子群（I，第74页）。证明同态 \(\hat{i}^q: \hat{\mathrm{H}}^q(\mathrm{G}, \mathrm{M}) \to \hat{\mathrm{H}}^q(\mathrm{H}, \mathrm{M})\) 的核是 \(l\) -准素分支（\(\mathbf{Z}\) -模 \(\hat{\mathrm{H}}^q(\mathrm{G}, \mathrm{M})\) 的， \(l \neq p\) ）之和。特别地，如果对每个素数 \(p\) 选择 G 的一个 \(p\) -Sylow子群 \(\mathrm{H}_p\) ，则映射 \(\hat{\mathrm{H}}^q(\mathrm{G}, \mathrm{M}) \to \prod \hat{\mathrm{H}}^q(\mathrm{H}_p, \mathrm{M})\) （由诸同态 \(\hat{i}^q\) 导出）是单射。
\end{enumerate}

\begin{enumerate}
\def\labelenumi{\arabic{enumi})}
\setcounter{enumi}{17}
\tightlist
\item
  设G是一个有限群，M是一个 \(\mathbf{Z}^{(G)}\) -模。称M是上同调平凡的，如果对G的每个子群H以及每个 \(q \in \mathbf{Z}\) ，群 \(\hat{\mathrm{H}}^q(\mathrm{H}, \mathrm{M})\) 为零。
\end{enumerate}

\begin{enumerate}
\def\labelenumi{\alph{enumi})}
\item
  证明一个 \(\mathbf{Z}^{(G)}\) -模相对于 \(\mathbf{Z}\) 是投射的（特别是诱导模，参见练习12）是上同调平凡的。
\item
  设n，r，k是三个整数，使得r整除 \((k^{n} - 1)\) ；设G是一个n阶循环群， \(\sigma\) 为 G 的一个生成元。将 Z-模 M = Z/rZ 赋予 \(\mathbf{Z}^{(G)}\) -模结构，它由 \(\sigma m = km\) （对所有 \(m \in M\) ）定义。证明：M 相对于 Z 是投射的，当且仅当 n 与 r 互素。
\item
  使用 b) 的记号，证明：为了使 \(\hat{\mathrm{H}}^{q}(\mathrm{G}, \mathrm{M}) = 0\) 对所有 \(q \in \mathbf{Z}\) 成立，其充分必要条件为 \((r, k - 1)(r, 1 + k + \cdots + k^{n-1}) = r\) 。特别地，如果 \(r = k^{n} - 1\) ，证明 M 在上同调意义下是平凡的。
\item
  由b)和c)，推出一个 \(\mathbf{Z}^{(G)}\) -模的例子，它是上同调平凡的，但相对于 \(\mathbf{Z}\) 不是投射的，以及一个 \(\mathbf{Z}^{(G)}\) -模M的例子，使得 \(\hat{\mathrm{H}}^q (\mathrm{G},\mathrm{M}) = 0\) 对所有 \(q\in \mathbf{Z}\) 但不是上同调平凡的。
\end{enumerate}

¶ 19) 设p为一个素数，G为一个p-群，M为一个Z \(^{(G)}\) -模。我们设 A = F \(_{p}^{(G)}\) .

\begin{enumerate}
\def\labelenumi{\alph{enumi})}
\tightlist
\item
  我们假设 \(pM = 0\) 。证明以下条件是等价的：
\end{enumerate}

α) 存在一个有理整数 q 使得 \(\hat{\mathrm{H}}^{q}(\mathrm{G},\mathrm{M})=0\) .

β) M 在上同调意义下是平凡的（练习 18）。

\(\gamma)\) M 是一个诱导 \(\mathbf{Z}^{(G)}\) -模。

δ) A-模 M 是自由的。

（为了证明 \(\alpha\) ）蕴含 \(\delta\) )，归结到 \(q = -2\) ；证明此时有

\[
\operatorname{Tor} _ {1} ^ {\mathbf {A}} \left(\mathbf {F} _ {p}, \mathbf {M}\right) = 0,
\]

并以第185页的练习4以及VIII章第8节的练习22得出结论。

\begin{enumerate}
\def\labelenumi{\alph{enumi})}
\setcounter{enumi}{1}
\tightlist
\item
  我们假设乘以 \(p\) 在 M 中是单射。证明以下条件是等价的：
\end{enumerate}

α) 存在整数 q 使得群 \(\hat{\mathrm{H}}^{q}(\mathrm{G},\mathrm{M})\) 和 \(\hat{\mathrm{H}}^{q+1}(\mathrm{G},\mathrm{M})\) 为零。

β) M 是上同调平凡的。

γ) A-模 M/pM 是自由的。

设 G 为有限群，M 为 \(\mathbf{Z}^{(G)}\) -模。

\begin{enumerate}
\def\labelenumi{\alph{enumi})}
\tightlist
\item
  我们假设 \(\mathbf{Z}\) -模 M 是自由的。证明： \(\mathbf{Z}^{(G)}\) -模 M 是投射的当且仅当对每个 Sylow 子群 H \(\subset\) G， \(\mathbf{Z}^{(H)}\) -模 M 都是上同调平凡的。
\end{enumerate}

（为证明该条件是充分的，考虑一个分解 \(0 \to R \to L \to M \to 0\) （\(\mathbf{Z}^{(G)}\) -模 M 的），其中 L 是自由的；证明 \(\mathbf{Z}^{(H)}\) -模 \(Q = \operatorname{Hom}_{\mathbf{Z}}(M, N)\) 是上同调平凡的，方法是观察到 \(Q/pQ\) 是上同调平凡的并利用习题 19。由习题 17 e) 推出 \(H^1(G, Q) = 0\) ，并由此得出 M 是投射的。）

\begin{enumerate}
\def\labelenumi{\alph{enumi})}
\setcounter{enumi}{1}
\tightlist
\item
  证明以下条件是等价的：
\end{enumerate}

\(\alpha)\) \(\mathbf{Z}^{(G)}\) -模 M 是上同调平凡的。

β) \(\mathbf{Z}^{(G)}\) -模 M 具有有限投射分解。

γ) \(\mathbf{Z}^{(G)}\) -模 M 具有长度为 1 的投射分解。

（为证明 \(\alpha\) ）蕴含 \(\gamma\) )，考虑正合序列 \(0 \to R \to L \to M \to 0\) ，其中 \(L\) 是自由 \(\mathbf{Z}^{(G)}\) -模，并将 \(a)\) 应用于 \(R\) 。）

\begin{enumerate}
\def\labelenumi{\alph{enumi})}
\setcounter{enumi}{2}
\item
  我们假设 Z-模 M 是可除的。证明： \(\mathbf{Z}^{(G)}\) -模 M 是上同调平凡的当且仅当它是内射的（考虑一个正合序列 \(0 \rightarrow M \rightarrow I \rightarrow Q \rightarrow 0\) ，其中 I 是内射 \(\mathbf{Z}^{(G)}\) -模，并利用 b) 证明 \(\mathbf{Z}^{(G)}\) -模 \(\operatorname{Hom}_{\mathbf{Z}}(\mathbf{Q}, \mathbf{M})\) 是上同调平凡的）。
\item
  一个 \(\mathbf{Z}^{(G)}\) -模是上同调平凡的，当且仅当它具有长度为 1 的内射分解。
\end{enumerate}

¶ 21) 设 G 为有限阶循环群。

\begin{enumerate}
\def\labelenumi{\alph{enumi})}
\item
  如果 M 是一个 \(\mathbf{Z}^{(\mathbf{G})}\) -模，且 \(q\) 是一个偶数（相应地，奇数）有理整数，则群 \(\hat{\mathrm{H}}^q (\mathbf{G},\mathbf{M})\) 同构于 \(\hat{\mathrm{H}}^0 (\mathbf{G},\mathbf{M})\) (相应地， \(\hat{\mathrm{H}}^1 (\mathbf{G},\mathbf{M}))\) （参见习题 7）。
\item
  我们用 \(\mathcal{C}\) 表示由这样的有限型 \(\mathbf{Z}^{(G)}\) -模 T 的类所构成的集合：群 \(\hat{\mathrm{H}}^q (\mathbf{G},\mathbf{T})\) 对每一个 \(q\in \mathbf{Z}\) 都是有限的；如果 \(\mathbf{Z}^{(G)}\) -模 M 是 \(\mathcal{C}\) 型的，我们设
\end{enumerate}

\[
h (\mathrm{M}) = \operatorname{Card} \left(\mathrm{H} ^ {0} (\mathrm{G}, \mathrm{M})\right) / \operatorname{Card} \left(\mathrm{H} ^ {1} (\mathrm{G}, \mathrm{M})\right).
\]

设 \(0 \to M' \to M \to M'' \to 0\) 为 \(\mathbf{Z}^{(G)}\) -模的一个正合序列。证明：如果模 \(M'\) , \(M, M''\) 中有两个属于类型 \(\mathcal{C}\) ，则第三个亦然，并且此时有 \(h(M) = h(M') h(M'')\) 。特别地，函数 \(h\) 定义了从 Grothendieck 群 \(\mathbf{K}(\mathcal{C})\) 到乘法群 \(\mathbf{Q}^*\) 的一个同态。c) 设 \(M\) 为一个有限 \(\mathbf{Z}^{(G)}\) -模。证明 \(M\) 是 \(\mathcal{C}\) 型的并且有 \(h(M) = 1\) 。

\begin{enumerate}
\def\labelenumi{\alph{enumi})}
\setcounter{enumi}{3}
\item
  假设 G 的阶是素数 p。我们用 \(\mathcal{C}'\) 表示由这样的有限生成 \(\mathbf{Z}^{(G)}\) -模的类所构成的集合：在其中乘以 \(p\) 的核和余核都是有限的；如果 M 是 \(\mathcal{C}'\) 型的，我们设 \(\varphi(M) = \text{Card (Coker } p_M)/\text{Card (Ker } p_M\) 。如 b) 中那样证明 \(\varphi\) 定义了从 \(\mathbf{K}(\mathcal{C}')\) 到 \(\mathbf{Q}^*\) 的一个同态。
\item
  证明 \(\mathbf{K}(\mathcal{C}')\) 由具有以下性质之一的 \(\mathbf{Z}^{(G)}\) -模 \(T\) 的类生成：
\end{enumerate}

\begin{enumerate}
\def\labelenumi{(\roman{enumi})}
\item
  T 是有限的。
\item
  乘以 \(p\) 在T中是双射的。
\item
  T = Z 赋予 G 的平凡作用。
\item
  \(\mathbf{T} = \mathbf{Z}^{(\mathbf{G})}\) .
\item
  \(\mathbf{T} = \mathbf{Q}_p / \mathbf{Z}_p\) （参见 TG, III, p. 84, 习题 23），赋予 G 的平凡作用。
\item
  \(\mathrm{T} = \mathbf{Z}^{(\mathrm{G})}\otimes_{\mathbf{Z}}(\mathbf{Q}_p / \mathbf{Z}_p)\) .
\end{enumerate}

（首先证明每个 \(\mathbf{Z}^{(\mathrm{G})}\) -模若为 \(\mathcal{C}'\) 型，则在 \(\mathbf{K}(\mathcal{C}')\) 中都可以写成一个 \(\mathbf{Z}^{(\mathrm{G})}\) -模（在 \(\mathbf{Z}\) 上有限生成）与一个 \(\mathbf{Z}^{(\mathrm{G})}\) -模（在其中乘以 \(p\) 是满射的）之和。然后确定环 \(\mathbf{Q}^{(\mathrm{G})}\) 和 \(\mathbf{Q}_p^{(\mathrm{G})}\) 上单模的结构。）

\begin{enumerate}
\def\labelenumi{\alph{enumi})}
\setcounter{enumi}{6}
\tightlist
\item
  从前述内容推出， \(\mathcal{C}' \subset \mathcal{C}\) ，并且如果 M 是一个 \(\mathbf{Z}^{(G)}\) -模且为 \(\mathcal{C}'\) 型，则有
\end{enumerate}

\[
h (\mathbf {M}) ^ {p - 1} = \left(\varphi \left(\mathrm{H} ^ {0} (\mathbf {G}, \mathbf {M})\right)\right) ^ {p} / \varphi (\mathbf {M}).
\]

（在 \(f\) ) 所考虑的六种情形中验证该公式。）

\begin{enumerate}
\def\labelenumi{\arabic{enumi})}
\setcounter{enumi}{21}
\tightlist
\item
  设 K 为交换域，L 为 K 的有限伽罗瓦扩张，G 为其伽罗瓦群。a) 证明 G-模的精确序列
\end{enumerate}

\[
0 \rightarrow \mu_ {n} (\mathrm{L}) \rightarrow \mathrm{L} ^ {*} \xrightarrow {x \mapsto x ^ {n}} \mathrm{L} ^ {* n} \rightarrow 0
\]

使得能够定义一个同构 \(k_{L}:(L^{n}\cap\mathbf{K}^{*})/\mathbf{K}^{*n}\to\mathbf{H}^{1}(\mathbf{G},\mu_{n}(\mathbf{L}))\) 。当 \(\mu_{n}(\mathbf{K})\) 有 \(n\) 个元素，并且将 \(\mathbf{H}^{1}(\mathbf{G},\mu_{n}(\mathbf{L}))\) 与 Hom \((\mathbf{G},\mu_{n}(\mathbf{K}))\) 等同时，证明 \(k_{L}\) 与 V，§ 11，第 8 号中定义的同态一致。

\begin{enumerate}
\def\labelenumi{\alph{enumi})}
\setcounter{enumi}{1}
\tightlist
\item
  设 \(p\) 为一个素数；假设 \(K\) 的特征为 \(p\) 。证明 G-模的正合序列（V，§ 11，第 9 号）
\end{enumerate}

\[
0 \to \mathbf {F} _ {p} \to \mathrm{L} \stackrel {{\mathfrak {p}}} {{\to}} \mathfrak {p} (\mathrm{L}) \to 0
\]

使得能够定义一个同构

\[
a _ {L}: (\mathfrak {p} (L) \cap K) / \mathfrak {p} (K) \rightarrow H ^ {1} (G, F _ {p})
\]

它与同上引用处中定义的同态一致，若将

\[
\mathbf {H} ^ {1} (\mathbf {G}, \mathbf {F} _ {p})
\]

与 Hom 等同

\[
(\mathbf {G}, \mathbf {F} _ {p})
\]

\begin{enumerate}
\def\labelenumi{\arabic{enumi})}
\setcounter{enumi}{22}
\tightlist
\item
  设G是一个群，M是一个以G为算子域的群（1，第29页，定义2）。我们用 \(^{g}m\) 表示 \(g \in G\) 与 \(m \in M\) 的复合。我们说M是一个G-群，如果有 \(^{h(g)m} = ^{hg}m\) （对 \(m \in M\) , \(g\) , \(h \in G\) ），换句话说，如果映射 \(\varphi\) （从 G 到 Aut(M)，由 \(\varphi(g), m = ^{g}m\) （对 \(g \in G\) , \(m \in M\) ）给出）是一个同态。
\end{enumerate}

如果 M 是一个 G-群，我们用 \(\mathbf{H}^{0}(\mathbf{G},\mathbf{M})\) 表示 M 的由元素 \(m\in M\) （满足 \(^{g}m=m\) ，对所有 \(g\in G\) ）组成的子群，并用 \(\mathbf{Z}^{1}(\mathbf{G},\mathbf{M})\) 表示从 G 到 M 的满足以下条件的映射 h 的集合：

\[
h (g g ^ {\prime}) = h (g) ^ {g} h (g ^ {\prime}) \quad \text { for } \quad g, g ^ {\prime} \in G.
\]

\begin{enumerate}
\def\labelenumi{\alph{enumi})}
\item
  我们说两个元素 \(h_1, h_2\) （\(Z^1(\mathbf{G}, \mathbf{M})\) 中的）是上同调的，如果存在一个元素 \(m\) （属于 \(\mathbf{M}\) ）使得 \(h_2(g) = m^{-1} h_1(g)^g m\) （对所有 \(g \in \mathbf{G}\) ）。证明这样定义的关系是 \(Z^1(\mathbf{G}, \mathbf{M})\) 上的一个等价关系；其商集记为 \(H^1(\mathbf{G}, \mathbf{M})\) 。\(H^1(\mathbf{G}, \mathbf{M})\) 中映射 \(i \in Z^1(\mathbf{G}, \mathbf{M})\) （它满足 \(i(g) = 1_{\mathbf{M}}\) ，对所有 \(g \in \mathbf{G}\) ）的类记为 \(e_{G,\mathbf{M}}\) 或简记为 \(e_{\mathbf{M}}\) 。
\item
  设 \(f: \mathbf{M} \to \mathbf{N}\) 为 G-群的一个同态（I，第30页，定义3）；证明 \(f\) 诱导一个群同态 \(f^{0}: \mathrm{H}^{0}(\mathbf{G}, \mathbf{M}) \to \mathrm{H}^{0}(\mathbf{G}, \mathbf{N})\) 以及一个映射 \(f^{1}: \mathrm{H}^{1}(\mathbf{G}, \mathbf{M}) \to \mathrm{H}^{1}(\mathbf{G}, \mathbf{N})\) 使得 \(f^{1}(e_{\mathbf{M}}) = e_{\mathbf{N}}\) .
\item
  设 \(\mathbf{M} \xrightarrow{i} \mathbf{N} \xrightarrow{p} \mathbf{Q}\) 为一个由 \(\mathbf{G}\) -群与同态构成的序列，其中 \(i\) 单射， \(p\) 满射，且
\end{enumerate}

\[
\operatorname{Im} (i) = \operatorname{Ker} (p).
\]

证明群 \(\mathbf{N}^{\prime}=p^{-1}(\mathbf{H}^{0}(\mathbf{G},\mathbf{Q}))\) 作用于 \(\mathbf{Z}^{1}(\mathbf{G},\mathbf{M})\) ，使得 \((n,h)(g)=n^{-1}h(g)^{g}n\) （对 \(n\in\mathbf{N}^{\prime}\) , \(h\in\mathbf{Z}^{1}(\mathbf{G},\mathbf{M})\) , \(g\in\mathbf{G}\) ）；证明该作用通过取商诱导出 \(\mathbf{H}^{0}(\mathbf{G},\mathbf{Q})\) 在 \(\mathbf{H}^{1}(\mathbf{G},\mathbf{M})\) 上的一个作用。映射 \(i^{1}:\mathbf{H}^{1}(\mathbf{G},\mathbf{M})\to\mathbf{H}^{1}(\mathbf{G},\mathbf{N})\) 通过取商定义了一个从 \(\mathbf{H}^{1}(\mathbf{G},\mathbf{M})/\mathbf{H}^{0}(\mathbf{G},\mathbf{Q})\) 到 \(\operatorname{Im}(i^{1})\) 上的双射。

\begin{enumerate}
\def\labelenumi{\alph{enumi})}
\setcounter{enumi}{3}
\tightlist
\item
  \(\mathbf{H}^0 (\mathbf{G},\mathbf{Q})\) 在 \(\mathbf{H}^{1}(\mathbf{G},\mathbf{M})\) 上的作用使我们可以定义一个映射 \(\partial :\mathbf{H}^{0}(\mathbf{G},\mathbf{Q})\to \mathbf{H}^{1}(\mathbf{G},\mathbf{M})\) ，它由 \(\partial (q) = q\cdot e_{\mathrm{M}}\) （对 \(q\in \mathbf{H}^0 (\mathbf{G},\mathbf{Q})\) ）给出。证明我们有
\end{enumerate}

\[
\operatorname{Im} \left(i ^ {0}\right) = \operatorname{Ker} \left(p ^ {0}\right), \operatorname{Im} \left(p ^ {0}\right) = \partial^ {- 1} (e _ {\mathrm{M}}), \quad \operatorname{Im} (\partial) = (i ^ {1}) ^ {- 1} \left(e _ {\mathrm{N}}\right), \quad \operatorname{Im} \left(i ^ {1}\right) = (p ^ {1}) ^ {- 1} \left(e _ {\mathrm{Q}}\right).
\]

\begin{enumerate}
\def\labelenumi{\alph{enumi})}
\setcounter{enumi}{4}
\item
  我们现在假设 \(i(\mathbf{M})\) 包含于 \(\mathbf{N}\) 的中心中。证明 \(\partial\) 是一个群同态；构造群 \(\mathrm{H}^{1}(\mathrm{G},\mathrm{M})\) 在集合 \(\mathrm{H}^{1}(\mathrm{G},\mathrm{N})\) 上的一个作用，使得映射 \(p^1\) 通过取商定义了一个从 \(\mathrm{H}^{1}(\mathrm{G},\mathrm{N}) / \mathrm{H}^{1}(\mathrm{G},\mathrm{M})\) 到 \(\operatorname{Im}(p^1)\) 上的双射。
\item
  设 E 为从 G 到 N 的映射 k 的集合，使得 \(p \circ k \in \mathbb{Z}^{1}(\mathbf{G}, \mathbf{Q})\) 。对于 \(k \in E\) ，证明存在一个元素 \(f_{k}\) 的 \(\mathbb{Z}^{2}(\mathbf{G}, \mathbf{M})\) 使得 \(i(f_{k}(x, y)) = k(x)^{x} k(y)(k(xy))^{-1}\) 当 \(x, y \in G\) 。证明映射 \(k \mapsto f_{k}\) 通过取商定义了一个映射 \(\partial^{1}: H^{1}(\mathbf{G}, \mathbf{Q}) \to H^{2}(\mathbf{G}, \mathbf{M})\) 。证明我们有 Im \((p^{1}) = (\partial^{1})^{-1}(0)\) 。
\end{enumerate}

\begin{enumerate}
\def\labelenumi{\arabic{enumi})}
\setcounter{enumi}{23}
\tightlist
\item
  设 K 是一个交换域，L 是 K 的有限伽罗瓦扩张，其伽罗瓦群为 G。
\end{enumerate}

设 V（相应地 V'）是一个 K-向量空间，t（相应地 t'）是 V（相应地 V'）上的一个 (p, q) 型张量（III，p.～63）。(V, t) 到 (V', t') 的一个 K-同构（或简称同构）是一个 K-线性同构 u : V → V'，使得 (T\^{}p(u) ⊗ T\^{}q(ü))(t) = t'。我们用 V\_L 表示 L-向量空间 L ⊗\_K V，并用 t\_L 表示 V\_L 上由 t 经标量扩张得到的 (p, q) 型张量。我们用 F\_\{L/K\}(V, t) 表示由满足 (W\_L, τ\_L) L-同构于 (V\_L, t\_L) 的偶 (W, τ) 的同构类所构成的集合。

\begin{enumerate}
\def\labelenumi{\alph{enumi})}
\tightlist
\item
  设 M 为 \((V_{L}, t_{L})\) 的 L-自同构群。对于 \(m \in M\) 和 \(\sigma \in G\) ，我们设
\end{enumerate}

\[
{ } ^ { \sigma } m = \pi _ { \mathrm{v} } ( \sigma ) \circ m \circ \pi _ { \mathrm{v} } ( \sigma ) ^ { - 1 } ,
\]

其中 \(\pi_{\mathbf{V}}: \mathbf{G} \to \operatorname{Aut}_{\mathbf{K}}(\mathbf{V}_{\mathbf{L}})\) 表示 \(\mathbf{G}\) 在 \(\mathbf{V}_{\mathbf{L}} = \mathbf{L} \otimes_{\mathbf{K}} \mathbf{V}\) 上的、由 \(\mathbf{G}\) 在 \(\mathbf{L}\) 上的作用导出的作用。证明这样就定义了 \(\mathbf{G}\) -群结构于 \(\mathbf{M}\) 上。

\begin{enumerate}
\def\labelenumi{\alph{enumi})}
\setcounter{enumi}{1}
\item
  设 (W, τ) 是 \(F_{L/K}(V, t)\) 的一个元素，并设 u 是从 \((\mathrm{V}_{\mathrm{L}}, t_{\mathrm{L}})\) 到 \((\mathrm{W}_{\mathrm{L}}, \tau_{\mathrm{L}})\) 上的一个 L-同构。对于 \(\sigma \in G\) ，我们用 \(f_{u}(\sigma)\) 表示 M 的元素 \(u^{-1} \circ \pi_{\mathbf{W}}(\sigma) \circ u \circ \pi_{\mathbf{V}}(\sigma)^{-1}\) 。证明 \(f_{u} \in Z^{1}(\mathbf{G}, \mathbf{M})\) （习题 23），并且它在 \(H^{1}(\mathbf{G}, \mathbf{M})\) 中的类不依赖于 u 的选择；把它记为 \(\theta(\mathbf{W}, \tau)\) 。
\item
  证明映射 \(\theta : F_{L/K}(V, t) \to H^1(G, M)\) 是双射的，并且有 \(\theta(V, t) = e_M\) .
\end{enumerate}

设 K 为一个交换域，L 为 K 的有限伽罗瓦扩张，G 为其伽罗瓦群。我们令 G 作用于群 GL(n, L)，规定 σ(aij) = (σ(aij))，其中 σ ∈ G，(aij) ∈ GL(n, L)。在此作用下稳定的 GL(n, L) 的每个子群因此具有 G-群结构（习题 23）。a) 证明集合 H¹(G, GL(n, L)) 与 H¹(G, SL(n, L)) 都归结为单个元素（参见 V，§ 10，第 5 号，命题 9）。

\begin{enumerate}
\def\labelenumi{\alph{enumi})}
\setcounter{enumi}{1}
\item
  证明集合 \(\mathbf{H}^{1}(\mathbf{G},\mathbf{PGL}(n,\mathbf{L}))\) 与中心单代数 S 的同构类集合等同，其中代数 \(\mathbf{L}\otimes_{\mathbf{K}}\mathbf{S}\) 同构于 \(\mathbf{M}_{n}(\mathbf{L})\) （利用习题 24）。通过将这样的代数与其在群 \(\mathbf{Br}(\mathbf{K},\mathbf{L})\) （VIII，§ 13）中的类对应，并通过将后一群与 \(\mathbf{H}^{2}(\mathbf{G},\mathbf{L}^{*})\) （同上）等同，得到映射 \(\mathbf{H}^{1}(\mathbf{G},\mathbf{PGL}(n,\mathbf{L}))\to \mathbf{H}^{2}(\mathbf{G},\mathbf{L}^{*})\) ，它是映射 \(\partial^{1}\) （习题 23， \(f\) ）的相反映射，后者由扩张 \(\mathbf{L}^{*}\to \mathbf{GL}(n,\mathbf{L})\to \mathbf{PGL}(n,\mathbf{L})\) 推导出。证明 \(\partial^{1}\) 的像由阶为 \(n\) 的元素（在 \(\mathbf{H}^{2}(\mathbf{G},\mathbf{L}^{*})\) 中）组成。
\item
  证明集合 \(\mathbf{H}^{1}(\mathbf{G},\mathbf{Sp}(2n,\mathbf{L}))\) 归结为单个元素（参见 IX，§ 5，第 1 号，定理 1）。
\item
  设 Q 是 K-向量空间 V 上的一个二次型， \(Q_{L}\) 是它在 \(L \otimes_{K} V\) 上导出的二次型。证明 \(\mathrm{H}^{1}(\mathrm{G}, \mathbf{O}(\mathrm{Q}_{\mathrm{L}}))\) 与二次型 \(Q'\) （V 上的，满足 \(Q_{L}'\) 同构于 \(Q_{L}\) ）的同构类集合等同。
\end{enumerate}

\begin{enumerate}
\def\labelenumi{\arabic{enumi})}
\setcounter{enumi}{25}
\tightlist
\item
  我们用 \(\mathbf{A}^e\) 表示 \(k\) -代数 \(\mathbf{A} \otimes_k \mathbf{A}^0\) ；每个 \((\mathbf{A}, \mathbf{A})\) -双模自然带有左 \(\mathbf{A}^e\) -模结构和右 \(\mathbf{A}^e\) -模结构。设 M 为一个 \((\mathbf{A}, \mathbf{A})\) -双模， \(n\) 为一个整数 \(\geqslant 0\) ；我们设：
\end{enumerate}

\[
\mathrm{H} _ {n} (\mathrm{A}, \mathrm{M}) = \operatorname{Tor} _ {n} ^ {\mathrm{A} ^ {e}} (\mathrm{M}, \mathrm{A}) \quad \mathrm{H} ^ {n} (\mathrm{A}, \mathrm{M}) = \operatorname{Ext} _ {\mathrm{A} ^ {e}} ^ {n} (\mathrm{A}, \mathrm{M}).
\]

\begin{enumerate}
\def\labelenumi{\alph{enumi})}
\item
  \(k\) -模 \(\mathbf{H}_0(\mathbf{A}, \mathbf{M})\) 同构于 \(\mathbf{M}\) 关于子 \(k\) -模（由元素 \((am - ma)\) （对 \(a \in \mathbf{A}, m \in \mathbf{M}\) ）生成）的商； \(k\) -模 \(\mathbf{H}^0(\mathbf{A}, \mathbf{M})\) 同构于子 \(k\) -模（\(\mathbf{M}\) 的，由元素 \(m\) （满足 \(am = ma\) ，对所有 \(a \in \mathbf{A}\) ）组成的）。\(k\) -模 \(\mathbf{H}^1(\mathbf{A}, \mathbf{M})\) 被等同于导子 \(k\) -模 \(\mathbf{D}_k(\mathbf{A}, \mathbf{M})\) 关于内导子所构成的子 \(k\) -模的商，内导子即那些导子 \(d_m\) （对 \(m \in \mathbf{M}\) ），定义为 \(d_m(a) = am - ma\) （对所有 \(a \in \mathbf{A}\) ）（参见III，第132页）。
\item
  假设A是一个投射k-模。证明标准分解B(A)（X，第58页）定义了左 \(\mathbf{A}^{e}\) 模A的一个投射分解。由此推出k-模H \(^{n}\) (A, M) 同构于复形 C(A, M) 的同调模，其中，对于 \(p \geqslant 0\) ，C \(^{p}\) (A, M) 是从 A \(^{p}\) 到 M 的 k-多重线性映射构成的 k-模，C \(^{p}\) (A, M) = 0 当 \(p < 0\) ，其微分由以下公式给出：
\end{enumerate}

当 \(f \in \mathbf{C}^p(\mathbf{A}, \mathbf{M})\) , \(x_0, \ldots, x_p \in \mathbf{A}\) .

给出 \(H_{n}(A, M)\) 的 k-模的类似表达式。

\begin{enumerate}
\def\labelenumi{\alph{enumi})}
\setcounter{enumi}{2}
\item
  在 \(b\) ) 的假设下，设 \(\rho : k \to k'\) 为交换环的一个同态，A' 为 \(k'\) -代数 \(k' \otimes_k \mathbf{A}\)', \(\mathbf{M}'\) 为一个 \((\mathbf{A}', \mathbf{A}')\) -双模。证明 \(k\) -模 \(\mathrm{H}_n(\mathbf{A}, \mathbf{M}')\) 和 \(\mathrm{H}_n(\mathbf{A}', \mathbf{M}')\) （相应地， \(\mathrm{H}^n(\mathbf{A}, \mathbf{M}')\) 和 \(\mathrm{H}^n(\mathbf{A}', \mathbf{M}')\) ）是同构的。
\item
  此后假设 k 是一个域。设 B 为第二个 k-代数（结合且含幺），M 为一个 (A, A)-双模，N 为一个 (B, B)-双模。对每一个 \(n \geqslant 0\) ，定义 k-模的同构
\end{enumerate}

\[
\mathrm{H} _ {n} (\mathrm{A} \otimes_ {k} \mathrm{B}, \mathrm{M} \otimes_ {k} \mathrm{N}) \rightarrow \bigoplus_ {p + q = n} \left(\mathrm{H} _ {p} (\mathrm{A}, \mathrm{M}) \otimes_ {k} \mathrm{H} _ {q} (\mathrm{B}, \mathrm{N})\right).
\]

若再假设 A 与 B 在 k 上是有限维的，证明同样存在一个 k-同构

\[
\mathrm{H} ^ {n} (\mathrm{A} \otimes_ {k} \mathrm{B}, \mathrm{M} \otimes_ {k} \mathrm{N}) \rightarrow \bigoplus_ {p + q = n} \left(\mathrm{H} ^ {p} (\mathrm{A}, \mathrm{M}) \otimes_ {k} \mathrm{H} ^ {q} (\mathrm{B}, \mathrm{N})\right) \quad \text { for   every } \quad n \geqslant 0
\]

（利用 X，第 76 页，定理 3）。

\begin{enumerate}
\def\labelenumi{\alph{enumi})}
\setcounter{enumi}{4}
\tightlist
\item
  设 M，N 为两个左 A-模，P 为右 A-模，使得 k-模 \(\operatorname{Hom}_{k}(\mathbf{M}, \mathbf{N})\) 和 \(M \otimes_{k} P\) 具有自然的 (A, A)-双模结构。对每个 \(n \geqslant 0\) 定义 k-模的同构：
\end{enumerate}

\[
\mathrm{H} ^ {n} (\mathrm{A}, \operatorname{Hom} _ {k} (\mathrm{M}, \mathrm{N})) \rightarrow \operatorname{Ext} _ {\mathrm{A}} ^ {n} (\mathrm{M}, \mathrm{N}) \quad \text { and } \quad \mathrm{H} _ {n} (\mathrm{A}, \mathrm{M} \otimes_ {k} \mathrm{P}) \rightarrow \operatorname{Tor} _ {n} ^ {\mathrm{A}} (\mathrm{P}, \mathrm{M}).
\]

¶ 27) 称 k-代数 A 的一个增广是 k-代数的一个含幺同态 π : A → k；称 (A, π) 为增广 k-代数。该增广赋予 k 一个左 A-模结构。设 M 为左 A-模，N 为右 A-模；令

\[
\mathrm{H} _ {n} (\mathrm{A}, \pi ; \mathrm{N}) = \operatorname{Tor} _ {n} ^ {\mathrm{A}} (\mathrm{N}, k) \quad \text { and } \quad \mathrm{H} ^ {n} (\mathrm{A}, \pi ; \mathrm{M}) = \operatorname{Ext} _ {\mathrm{A}} ^ {n} (k, \mathrm{M}) \quad \text { for   every } \quad n \geqslant 0.
\]

\begin{enumerate}
\def\labelenumi{\alph{enumi})}
\tightlist
\item
  我们用 \(\mathbf{M}_{(\pi)}\) (相应地， \(\mathbf{N}_{(\pi)}\) ) 表示群 \(\mathbf{M}\) (相应地， \(\mathbf{N}\) )，赋予它由下式定义的 (A, A)-双模结构：
\end{enumerate}

\[
a m a ^ {\prime} = \pi (a ^ {\prime}) a m (\text { resp. } a n a ^ {\prime} = n \pi (a) a ^ {\prime} \quad \text { for } \quad a, \quad a ^ {\prime} \in \mathbf {A}, \quad m \in \mathbf {M}, \quad n \in \mathbf {N}).
\]

证明 \(k\) -模 \(\mathbf{H}_0(\mathbf{A}, \pi; \mathbf{N})\) (相应地， \(\mathbf{H}^0(\mathbf{A}, \pi; \mathbf{M})\) ，相应地 \(\mathbf{H}^1(\mathbf{A}, \pi; \mathbf{M})\) )同构于练习26中定义的 \(k\) -模 \(\mathbf{H}_0(\mathbf{A}, \mathbf{N}_{(\pi)})\) (相应地， \(\mathbf{H}^0(\mathbf{A}, \mathbf{M}_{(\pi)})\) ，相应地 \(\mathbf{H}^1(\mathbf{A}, \mathbf{M}_{(\pi)})\) )。 \(b)\) 我们现在假设 \(\mathbf{A}\) 是一个投射 \(k\) -模。使用标准分解 (X, p.～58) 定义 \(k\) -同构 \(\varphi_n: \mathbf{H}_n(\mathbf{A}, \mathbf{N}_{(\pi)}) \to \mathbf{H}_n(\mathbf{A}, \pi; \mathbf{N})\) 和 \(\varphi^n: \mathbf{H}^n(\mathbf{A}, \mathbf{M}_{(\pi)}) \to \mathbf{H}^n(\mathbf{A}, \pi; \mathbf{M})\) （对所有 \(n \geqslant 0\) ）。

\begin{enumerate}
\def\labelenumi{\alph{enumi})}
\setcounter{enumi}{2}
\item
  我们假设存在一个 \(k\) -代数的同态 \(\rho: \mathbf{A} \to \mathbf{A}^e\) ，使得 \(\mathbf{A}^e\) 是一个投射右 \(\mathbf{A}\) -模，且使得 \(\mathbf{A}^e\) 中由 \(\rho(\text{Ker } \pi)\) 生成的左理想等于同态 \(\mathbf{A} \otimes_k \mathbf{A}^* \to \mathbf{A}\) （由 \(\mathbf{A}\) 中的乘法所定义）的核。证明对于每一个 \((\mathbf{A}, \mathbf{A})\) -双模 \(\mathbf{Q}\) 及每个整数 \(n \geqslant 0\) 有 \(k\) -同构 \(\mathrm{H}_n(\mathbf{A}, \pi; \mathbf{Q}) \to \mathrm{H}_n(\mathbf{A}, \mathbf{Q})\) (相应地， \(\mathrm{H}^n(\mathbf{A}, \pi; \mathbf{Q}) \to \mathrm{H}^n(\mathbf{A}, \mathbf{Q})\) )，这里 \(\mathbf{Q}\) 被视为右侧（相应地，左侧）的 \(\mathbf{A}\) -模，且经由同态 \(\rho\) 。（使用 \(\mathbf{X}\) ，第109页的典范同态。）
\item
  证明以下情形满足c)的假设条件：
\end{enumerate}

\(\alpha)\) A 是（在 \(k\) 上的）群 \(\mathbf{G}\) 的代数， \(\rho\) 定义为 \(\rho(e_g) = e_g \otimes e_{g-1}\) 。

β) A 是投射 k-模 V 的张量（相应地，对称）代数；ρ 由下式定义

\[
\rho (v) = v \otimes 1 - 1 \otimes v.
\]

* γ) A 是 k 上李代数 g 的包络代数，作为 k-模是自由的；我们有

\[
\rho (x) = x \otimes 1 - 1 \otimes x \quad \text { for } \quad x \in \mathfrak {g}
\]

(利用LIE I，§2，第7号中定理1的推论5)。

δ) A 是域 k 上光滑代数群 G 的函数环；我们取

\[
\rho (f) = \sum_ {i = 1} ^ {k} f _ {i} ^ {\prime} \otimes f _ {i} ^ {\prime \prime}, \quad \text { where } \quad f (g h ^ {- 1}) = \sum_ {i = 1} ^ {k} f _ {i} ^ {\prime} (g) f _ {i} ^ {\prime \prime} (h) \quad \text { for } \quad g, h \in G. *
\]

设g是交换环k上的李代数，U是其包络代数。我们赋予k以对应于g的平凡表示的左U-模结构。设N为右g-模，M为左g-模，n为整数≥0；我们令

\[
\mathrm{H} _ {n} (\mathfrak {g}, \mathrm{N}) = \operatorname{Tor} _ {n} ^ {\mathrm{U}} (\mathrm{N}, k) \quad \text { and } \quad \mathrm{H} ^ {n} (\mathfrak {g}, \mathrm{M}) = \operatorname{Ext} _ {\mathrm{U}} ^ {n} (k, \mathrm{M}).
\]

\begin{enumerate}
\def\labelenumi{\alph{enumi})}
\item
  \(k\) -模 \(\mathbf{H}_0(\mathfrak{g}, \mathbf{N})\) (相应地， \(\mathbf{H}^0(\mathfrak{g}, \mathbf{M})\) )同构于 \(\mathbf{N}/\mathbf{Ng}\) （相应地，\(\mathbf{M}\) 的由被 \(\mathfrak{g}\) 所零化的元素组成的子模）。我们用 \(\mathbf{Z}^1(\mathfrak{g}, \mathbf{M})\) (相应地， \(\mathbf{B}^1(\mathfrak{g}, \mathbf{M})\) ) 表示子 \(k\) -模（\(\mathrm{Hom}_k(\mathfrak{g}, \mathbf{M})\) 的），它由元素 \(f\) （满足 \(f([x, y]) = xf(y) - yf(x)\) ，对 \(x, y \in \mathfrak{g}\) ）组成（相应地，由元素 \(f_m\) （对 \(m \in \mathbf{M}\) ，满足 \(f_m(x) = xm\) ，对 \(x \in \mathfrak{g}\) ）组成）。证明 \(\mathbf{H}^1(\mathfrak{g}, \mathbf{M})\) 与商 \(\mathbf{Z}^1(\mathfrak{g}, \mathbf{M})/\mathbf{B}^1(\mathfrak{g}, \mathbf{M})\) 等同。
\item
  我们现在假设 \(\mathfrak{g}\) 是一个自由 \(k\) -模。证明 \(k\) -模 \(\mathbf{H}^n(\mathfrak{g}, \mathbf{M})\) 与复形 \(\mathbf{C}(\mathfrak{g}, \mathbf{M})\) 的同调模等同，其中对于 \(p \geqslant 0\) , \(\mathbf{C}^p(\mathfrak{g}, \mathbf{M})\) 是这样的 \(k\) -模：它由从 \(\mathfrak{g}^p\) 到 \(\mathbf{M}\) 的交替多重线性映射组成， \(\mathbf{C}^p(\mathfrak{g}, \mathbf{M}) = 0\) 当 \(p < 0\) ，其微分由以下公式给出
\end{enumerate}

\[
\begin{array}{r l} d f (x _ {1}, \dots , x _ {p + 1}) = & \sum_ {1 \leqslant i \leqslant p + 1} (- 1) ^ {i + 1} x _ {i}. f (x _ {1}, \dots , \hat {x} _ {i}, \dots , x _ {p + 1}) \\ & + \sum_ {1 \leqslant i <   j \leqslant p + 1} (- 1) ^ {i + j} f ([ x _ {i}, x _ {j} ], x _ {1}, \dots , \hat {x} _ {i}, \dots , \hat {x} _ {j}, \dots , x _ {p + 1}) \end{array}
\]

当 \(f \in C^{p}(\mathfrak{g}, \mathbf{M}), x_{1}, \ldots, x_{p+1} \in \mathfrak{g}\) （其中字母上的符号表示该字母应被省略）。

给出 \(k\) -模 \(\mathbf{H}_p(\mathfrak{g}, \mathbf{M})\) 的类似表达式。（使用 \(\mathbf{X}\) ，第184页，练习21。） c) 我们假设 \(\dim_k(\mathfrak{g}) = n < \infty\) 。证明 \(\mathbf{H}^p(\mathfrak{g}, \mathbf{M}) = 0\) 对所有 \(p \geqslant n + 1\) 且对每个 \(\mathfrak{g}\) -模 \(\mathbf{M}\) 成立，且存在一个 \(\mathfrak{g}\) -模 \(\mathbf{N}\) 使得 \(\mathbf{H}^n(\mathfrak{g}, \mathbf{N}) \neq 0\) （取 \(\mathbf{N} = \Lambda^n\mathfrak{g}\) ，其中 \(\mathfrak{g}\) 通过公式

\[
x _ {*} \left(x _ {1} \wedge \dots \wedge x _ {n}\right) = \sum_ {i = 1} ^ {n} x _ {1} \wedge \dots \wedge x _ {i - 1} \wedge [ x, x _ {i} ] \wedge x _ {i + 1} \wedge \dots \wedge x _ {n};
\]

作用；证明 \(df = 0\) 对所有 \(f \in \mathbf{C}^{n-1}(\mathfrak{g}, \mathbf{N})\) 。

\begin{enumerate}
\def\labelenumi{\alph{enumi})}
\setcounter{enumi}{3}
\tightlist
\item
  g 的伴随表示以及 g 在 M 上的表示定义了 g 在 C(g, M) 中的一个表示 θ，使得
\end{enumerate}

\[
\left(\theta (x) f\right) \left(x _ {1}, \dots , x _ {p}\right) = x. f (x _ {1}, \dots , x _ {p}) - \sum_ {i = 1} ^ {p} f (x _ {1}, \dots , [ x, x _ {i} ], \dots , x _ {p})
\]

当 \(f \in \mathbf{C}^p(\mathfrak{g}, \mathbf{M})\) , \(x, x_1, \ldots, x_p \in \mathfrak{g}\) 。证明：对每个 \(x \in \mathfrak{g}\) ， \(\theta(x)\) 是复形 \(\mathbf{C}(\mathfrak{g}, \mathbf{M})\) 的一个同伦于零的自同态（定义一个同伦 \(i(x)\) ，令 \((i(x) f) (x_1, \ldots, x_p) = f(x, x_1, \ldots, x_p)\) ）。 * 29) 我们沿用前一练习的记号；此外假设 \(k\) 是特征零的域，且李代数 \(\mathfrak{g}\) 是半单的（LIE, I, § 6）。设 M 是 \(k\) 上有限维的 g-模。

\begin{enumerate}
\def\labelenumi{\alph{enumi})}
\item
  若 M 是单的且非平凡的，则 \(H^{p}(g, M)=0\) 对每个 p\textgreater0 成立（与 M 相关联的 g 上的双线性形式是非退化的（LIE, I, § 6, n° 1, 命题 1）；注意乘以相应的 Casimir 元素（LIE, I, § 3, n° 7）在 k 中为零而在 M 上是双射的）。
\item
  证明 \(\mathbf{H}^{1}(\mathfrak{g},\mathbf{M}) = 0\) 对每个 \(\mathbf{M}\) 成立（利用 \(a\) ）以及 \(\mathfrak{g} = \mathcal{D}\mathfrak{g}\) 这一事实。反之，一个 \(k\) -李代数 \(\mathfrak{h}\) （有限维的），若 \(\mathbf{H}^{1}(\mathfrak{h},\mathbf{M}) = 0\) 对所有有限维 \(\mathfrak{h}\) -模 \(\mathbf{M}\) 成立，则它是半单的。
\item
  我们赋予 k 对应于平凡表示的 g-模结构。证明：对每个 \(p \geqslant 0\) ，k-向量空间 \(H^{p}(g, k)\) 同构于 \(C^{p}(g, k)\) 中由不变的交错 p-线性形式 f 组成的子空间，即满足
\end{enumerate}

\[
\sum_ {i = 1} ^ {p} f (x _ {1}, \dots , [ x, x _ {i} ], \dots , x _ {p}) = 0 \quad \text { for   all } \quad x, x _ {1}, \dots , x _ {p} \in \mathfrak {g}.
\]

的形式。（由 g 的表示的半单性以及练习 28, d) 推出 \(\mathbf{H}^p(\mathfrak{g}, k)\) 同构于 \(\mathbf{H}^p(\mathbf{C}_{\theta}(\mathfrak{g}, k))\) ，其中 \(\mathbf{C}_{\theta}(\mathfrak{g}, k)\) 是 \(\mathbf{C}(\mathfrak{g}, k)\) 中被所有自同态 \(\theta(x)\) （对 \(x \in \mathfrak{g}\) ）零化的子复形；然后证明 \(\mathbf{C}_{\theta}(\mathfrak{g}, k)\) 的微分为零。）

\begin{enumerate}
\def\labelenumi{\alph{enumi})}
\setcounter{enumi}{3}
\tightlist
\item
  由 c) 推出：对 k 上每个有限维 g-模 M， \(\mathrm{H}^{2}(\mathrm{g},\mathrm{M})=0\) 。（证明 g 上每个不变双线性形式都是对称的，办法是归结到 g 为单且 k 为代数闭的情形。） *
\end{enumerate}
% semantic-fragments: legacy-input-seam
\relax{}
